\section{Positioning and Design Principle}
\label{sec:positioning}

\subsection{Continuity with earlier TLRA}

The earlier TLRA line established three principles that remain central here. First, hallucination mitigation should act \textbf{during generation}, not only through post-hoc filtering. Second, the relevant unit of control is \textbf{candidate-local}, because hallucination often emerges at the level of specific lexical alternatives rather than global sequence quality. Third, any practical decode-time intervention must preserve linguistic structure rather than indiscriminately suppress visually weak tokens.

UniGround preserves these principles while changing the location and role of the learned component. It is formulated as a decode-time grounding-control interface rather than as a theory of how evidence should be rerouted inside the host MLLM. This distinction is substantive. The goal is not to reinterpret the host model's internal reasoning trajectory, but to attach a portable external controller that can evaluate and modulate candidate spans under a strict universality constraint.

\subsection{Why internal lexical alignment is not universal}

A learned module fails to be strictly universal if either its inputs or outputs are host-specific. More precisely, universality is broken when:
\begin{enumerate}
  \item the module consumes hidden states whose semantics depend on the host architecture; or
  \item the module emits vectors interpreted directly in the host model's lexical space.
\end{enumerate}

The point can be formalized directly. Let $m$ denote a host MLLM with hidden-state space $\mathcal{H}_m$, vocabulary $\mathcal{V}_m$, and unembedding map $U_m$. Suppose a learned controller $f_\theta$ consumes $h_t \in \mathcal{H}_m$ and emits a correction in the host lexical space $\mathbb{R}^{|\mathcal{V}_m|}$. If there is no model-independent isomorphism that simultaneously aligns $\mathcal{H}_m$, $\mathcal{H}_{m'}$, $\mathcal{V}_m$, and $\mathcal{V}_{m'}$ across two hosts $m$ and $m'$, then the same parameters $\theta$ cannot preserve semantic equivalence across both models. The learned mapping is therefore host-coupled by construction.

\texttt{Phi\_calib} instantiates exactly this host-coupled pattern. It reads model-native multimodal hidden states and projects them into the active model's lexical geometry. Its behavior therefore depends on hidden-state provenance, vocabulary structure, tokenizer conventions, and \texttt{lm\_head} organization. It remains an informative internal control, but it cannot ground a strict universal-plugin claim.

\subsection{Universality through standardized observation and action spaces}

To make universality defensible, the learned plugin must instead be defined through model-agnostic observation and action spaces:
\begin{itemize}
  \item \textbf{Observation space:} pixels, frozen image embeddings, frozen text embeddings, decoded strings, and optional frozen region proposals.
  \item \textbf{Action space:} span-level estimates of visual support, contradiction, and abstention.
\end{itemize}

Under this design, the learned plugin never consumes model-native hidden states, never reads \texttt{lm\_head.weight}, and never relies on a learned model-specific bridge. Universality is thus enforced jointly at the level of representation and attachment. The method is universal not because every component is identical across models, but because all model-dependent variation is confined to deterministic, non-learned attachment logic.
